\documentclass[12pt, reqno]{amsart}	
\usepackage{latexsym, amssymb, amscd, xypic, amsthm, setspace}
\input xy
\xyoption{all}

%%: Font = times
\usepackage{mathptmx} % rm & math
\usepackage[scaled=0.90]{helvet} % ss
\usepackage{courier} % tt
\normalfont
\usepackage[T1]{fontenc}
\pagestyle{plain}


\newtheorem{theorem}{Theorem}[subsection]
\newtheorem{prop}[theorem]{Proposition}
\newtheorem{lemma}[theorem]{Lemma}
\newtheorem{corollary}[theorem]{Corollary}
\newtheorem{claim}[theorem]{Claim}

\theoremstyle{definition}
\newtheorem{de}[theorem]{Definition}
\newtheorem{notation}{Notation}
\newtheorem{example}[theorem]{Example}
\newtheorem{examples}{Examples}
\newtheorem{remark}[theorem]{Remark}
\newtheorem{remarks}{Remarks}
\newtheorem{question}{Question}


\newtheorem*{nthm}{Theorem}
\newtheorem*{nlem}{Lemma}
\newtheorem*{nprop}{Proposition}
\newtheorem*{ncor}{Corollary}
\newtheorem*{nclaim}{Claim}
\newtheorem*{nde}{Definition}
\newtheorem*{nex}{Example}

\newtheoremstyle{test}{3pt}{3pt}{\itshape}{\parindent}{\bfseries}{---}{  }{\thmnote{#3}} \theoremstyle{test}
\newtheorem*{mythm}{}


%: paragraphs

\newcommand{\point}{\vspace{3mm}\par \noindent (\refstepcounter{subsection}{\bf \thesubsection) } }

\newcommand{\tpoint}[1]{\vspace{3mm}\par \noindent \refstepcounter{subsubsection}{\bf \thesubsubsection.} 
  {\em #1. ---} }
\newcommand{\epoint}[1]{\vspace{3mm}\par \noindent \refstepcounter{subsection}{\bf \thesubsection.} 
  {\em #1.} }
\newcommand{\spoint}{\vspace{3mm}\par \noindent \refstepcounter{subsection}{\bf \thesubsection.} }

\newcommand{\bpoint}[1]{\vspace{3mm}\par \noindent \refstepcounter{subsection}{\bf \thesubsection.} 
  {\bf #1.} }



%: numbering and margins

\numberwithin{equation}{subsection}

\setlength{\oddsidemargin}{0cm} \setlength{\evensidemargin}{0cm}
\setlength{\marginparwidth}{0in}
\setlength{\marginparsep}{0in}
\setlength{\marginparpush}{0in}
\setlength{\topmargin}{0in}
\setlength{\headheight}{0pt}
\setlength{\headsep}{0pt}
\setlength{\footskip}{.3in}
\setlength{\textheight}{9.2in}
\setlength{\textwidth}{6.5in}
\setlength{\parskip}{0.25pt}
\setlength{\parindent}{0.25in}

%: math shortcuts 


\DeclareMathOperator{\inv}{inv}
\DeclareMathOperator{\out}{Out}
\DeclareMathOperator{\gal}{Gal}
\renewcommand{\mod}{\mathrm{\, mod \, }}
\renewcommand{\d}[1]{{}^{#1}}
\renewcommand{\b}[1]{\mathbf{{#1}}}
\newcommand{\C}{\mathbb{C}}
\newcommand{\be}[1]{\begin{eqnarray} \label{#1}}
\newcommand{\ee}{\end{eqnarray}}
\newcommand{\re}{\mathsf{Re}}
\newcommand{\im}{\mathsf{Im}} 
\newcommand{\rr}{\rightarrow}

\newcommand{\n}{\mathbb{N}}
\newcommand{\zee}{\mathbb{Z}}
\newcommand{\R}{\mathbb{R}}
\newcommand{\pr}{\mathbb{P}}
\renewcommand{\gcd}[2]{\mathsf{gcd}(#1, #2)}
\newcommand{\ov}[1]{\overline{#1}}
\renewcommand{\Re}{\mathrm{Re}}
\renewcommand{\Im}{\mathrm{Im}}



\title{Problem Set 1 \\ MATH 411: Honors Complex Variables, Fall 2026}
\date{last compiled: \today}

\numberwithin{equation}{section}
\refstepcounter{section}

\begin{document}						%%
\maketitle



\begin{enumerate}

\item
\begin{enumerate} 

\item Put $-1 - i$ into polar form
\item Put $e^{- 5 i \pi /4}$ into the form $x + i y$.
\item Find the real and imaginary parts of $i^{1/4}$, taking the fourth root whose angle lies between $0$ and $\pi/2$.
\item Describe geometrically the set of points $z \in \mathbb{C}$ such that $|z| = \Re(z)+1$.
\end{enumerate} 
\item Let $f(z) = e^z$. Describe the image under $f$ of the following sets 
\begin{enumerate}
\item The set of $z = x + i y $ such that $x \leq 1$ and $0 \leq y \leq \pi$
\item The set of $z = x + i y $ such that $0 \leq y \leq \pi$ (no condition on $x$)
\end{enumerate}

\item Let $f$ be the function defined by $$ f(z) = \lim_{ n \rr \infty} \frac{1}{1 + n^2 z} .$$
 Show that $f$ is the characteristic function of the set $\{ 0 \}$, i.e. $f(0)=1$ and $f(z)=0$ otherwise. 
 
 \item \begin{enumerate}
\item Let $z, w$ be two complex numbers such that $\overline{z} w \neq 1$. Prove that $$ \left| \frac{w -z }{1 - \overline{w} z} \right| < 1 \text{ if } |z| < 1 \text{ and } |w| < 1 $$ and also that $$ \left| \frac{w -z }{1 - \overline{w} z} \right| = 1 \text{ if } |z| = 1 \text{ or } |w| = 1 $$ 

\emph{Hint: reduce to the case where $z$ is real, say $z=r$, in which case it suffices to prove that $$ (r - w) (r - \overline{w}) \leq ( 1- rw) (1 - r \overline{w}) $$ with equality when $|r|=1$ or $|w|=1$.}

\item Prove that for a fixed $w$ in the unit disc $\mathbb{D}$ the mapping $$ F: z \mapsto \frac{w - z}{1 - \overline{w} z} $$ satisfies the following conditions: 

\begin{enumerate} 
\item $F$ maps $\mathbb{D}$ to itself and is holomorphic on $\mathbb{D}$
\item $F$ interchanges $0$ and $w$ (i.e. $F(0)=w$ and $F(w)=0$)
\item $|F(z)|=1$ if $|z|=1$
\item $F: \mathbb{D} \rr \mathbb{D}$ is bijective (Hint: calculate $F \circ F$).
\end{enumerate}


\end{enumerate}

\item Show that in polar coordinates, i.e. if $F(x, y)$ is regarded as $F(r, \theta)= ( u(r,\theta), v(r, \theta))$, the Cauchy--Riemann equations take the form  $$ \frac{\partial u}{\partial r}  = \frac{1}{r} \frac{ \partial v}{\partial \theta} \text{ and } \frac{1}{r} \frac{ \partial u}{\partial \theta} = - \frac{\partial v}{\partial r}.$$  Use these equations to show that the logarithm function defined as $$ \log(z) := \log (r) + i \theta \text{ where } z = r e^{ i \theta} \text{ with } - \pi < \theta < \pi $$ is holomorphic in the region $r > 0$ and $- \pi < \theta < \pi$. 

\item \begin{enumerate}
\item Show that $$4 \frac{\partial}{\partial z} \frac{\partial}{\partial \overline{z}} = 4 \frac{\partial}{\partial \overline{z}} \frac{\partial}{\partial z} = \Delta $$ where $\Delta$ is the Laplacian, $$ \Delta = \frac{\partial^2}{\partial x^2} + \frac{\partial^2}{\partial y^2} .$$ \item Show that if $f$ is holomorphic on an open set $U \subset \mathbb{C}$, then $\psi(x, y) = \Re(f)$ and $\phi(x, y) = \Im(f)$ are both harmonic functions (i.e. they satisfy $\Delta \psi = \Delta \phi = 0$).
\end{enumerate}

\item Show that $f(x+ i y) = \sqrt{|x| | y |}$ satisfies the Cauchy-Riemann equations at the origin but that $f$ is not holomorphic at $0$

\item Suppose that $f$ is holomorphic in an open disc $\mathcal{D} \subset \mathbb{C}$, say $f= u(x, y) +  i v(x, y)$. Prove that \begin{enumerate} 
\item if $|f|^2 = u^2 + v^2$ is constant, then $f$ is constant;
\item if there exist real, non-zero constants $a, b,$ and $c$ such that $$ a u (x, y) + b v (x, y) = c, $$ then $f$ is constant. 
\end{enumerate}

\end{enumerate}


\end{document}
